\documentclass[conference,a4paper,10pt]{IEEEtran}
\IEEEoverridecommandlockouts

\usepackage[T1]{fontenc}
\usepackage[utf8]{inputenc}
\usepackage[brazil]{babel}
\usepackage{graphicx}
\usepackage{array}
\usepackage{booktabs}
\usepackage{multirow}
\usepackage{pdflscape}
\usepackage{amsmath}
\usepackage{xcolor}
\usepackage{listings}
\usepackage{url}
\usepackage[hidelinks]{hyperref}

\newcolumntype{L}[1]{>{\raggedright\arraybackslash}p{#1}}
\newcolumntype{C}[1]{>{\centering\arraybackslash}p{#1}}

\lstset{
  basicstyle=\ttfamily\scriptsize,
  breaklines=true,
  columns=fullflexible,
  frame=single,
  framesep=3pt,
  keywordstyle=\bfseries,
  language=SQL,
  showstringspaces=false,
  captionpos=b
}

\renewcommand{\IEEEkeywordsname}{Palavras-chave}

\begin{document}

\title{Projeto de Banco de Dados para uma Locadora de Ve\'iculos:\\ Dicion\'ario de Dados, Modelagem Entidade-Relacionamento e Normaliza\c{c}\~ao}

\author{
\IEEEauthorblockN{Marcello Azevedo Pinheiro Siqueira\thanks{Grupo constitu\'ido por tr\^es integrantes mediante exce\c{c}\~ao autorizada pelo docente em raz\~ao do n\'umero de alunos da turma.}}
\IEEEauthorblockA{Matr\'icula 24101166\\
Engenharia de Software\\
Instituto Brasileiro de Ensino,\\ Desenvolvimento e Pesquisa (IDP)\\
Bras\'ilia, DF, Brasil}
\and
\IEEEauthorblockN{Lucas Basile}
\IEEEauthorblockA{Matr\'icula 24101316\\
Engenharia de Software\\
Instituto Brasileiro de Ensino,\\ Desenvolvimento e Pesquisa (IDP)\\
Bras\'ilia, DF, Brasil}
\and
\IEEEauthorblockN{Miguel Matos}
\IEEEauthorblockA{Matr\'icula 24201146\\
Engenharia de Software\\
Instituto Brasileiro de Ensino,\\ Desenvolvimento e Pesquisa (IDP)\\
Bras\'ilia, DF, Brasil}
}

\maketitle

\begin{abstract}
Este trabalho apresenta o projeto do banco de dados de uma locadora de ve\'iculos com atua\c{c}\~ao em m\'ultiplas filiais, desenvolvido como Laborat\'orio 01 da disciplina de Banco de Dados. O projeto cobre tr\^es etapas. A primeira \'e o levantamento das entidades do cen\'ario e a constru\c{c}\~ao do dicion\'ario de dados, com fichas de campos, chaves, relacionamentos, regras de neg\'ocio, dom\'inios controlados e \'indices. A segunda \'e a modelagem entidade-relacionamento, com sete entidades e cardinalidades 1:1, 1:N e N:N explicitadas. A terceira \'e a normaliza\c{c}\~ao, partindo de uma estrutura \'unica n\~ao normalizada e aplicando a Primeira, a Segunda e a Terceira Formas Normais, com justificativa t\'ecnica de cada decomposi\c{c}\~ao. O modelo final possui sete rela\c{c}\~oes em 3FN, atende \`as quatro restri\c{c}\~oes funcionais definidas no enunciado e serve de base para o script de cria\c{c}\~ao em MySQL 8.
\end{abstract}

\begin{IEEEkeywords}
banco de dados, dicion\'ario de dados, modelo entidade-relacionamento, normaliza\c{c}\~ao, formas normais, MySQL
\end{IEEEkeywords}
